\documentclass[11pt]{article}
\usepackage[a4paper,margin=0.9in]{geometry}
\usepackage{amsmath,amssymb,amsthm,mathtools}
\usepackage[dvipsnames]{xcolor}
\usepackage{enumitem}
\usepackage{booktabs}
\usepackage{tcolorbox}
\usepackage{fancyhdr}
\usepackage[colorlinks=true,linkcolor=Blue]{hyperref}
\tcbuselibrary{skins,breakable}

\setlist{itemsep=3pt,topsep=4pt}
\pagestyle{fancy}\fancyhf{}
\lhead{\small Bayes Estimation --- Crash Course}
\rhead{\small Parametric Inference, B.Stat.\ 3rd yr}
\cfoot{\small\thepage}

\newcommand{\E}{\mathbb{E}}
\newcommand{\Var}{\operatorname{Var}}
\newcommand{\Prob}{\mathbb{P}}
\newcommand{\R}{\mathbb{R}}
\newcommand{\Xbar}{\overline{X}}
\newcommand{\iid}{\overset{\text{iid}}{\sim}}
\newcommand{\I}{\mathbb{I}}

\newtcolorbox{idea}[1]{colback=Blue!4,colframe=Blue!55,title={#1},
  fonttitle=\bfseries,breakable}
\newtcolorbox{confusion}[1]{colback=BrickRed!5,colframe=BrickRed!60,
  title={Common confusion: #1},fonttitle=\bfseries,breakable}
\newtcolorbox{recipe}[1]{colback=ForestGreen!5,colframe=ForestGreen!55,title={#1},
  fonttitle=\bfseries,breakable}

\begin{document}

\begin{center}
{\LARGE\bfseries Bayes Estimation}\\[3pt]
{\large A Crash Course from First Principles}\\[5pt]
{\small Parametric Inference, B.Stat.\ 3rd Year --- Lectures 6--7}
\end{center}

\vspace{4pt}\hrule\vspace{10pt}

\noindent
This assumes you know nothing about Bayes methods and everything about the first half of the
course (risk, unbiasedness, sufficiency, UMVUE). Read it front to back --- each section is
built on the one before. Worked examples use real numbers so you can check every step.

\tableofcontents

%%%%%%%%%%%%%%%%%%%%%%%%%%%%%%%%%%%%%%%%%%%%%%%%%%%%%%%%%
\section{Why Bayes exists at all}

Go back to the very first thing proved in Lecture 1.

\begin{idea}{The problem the whole course is trying to solve}
For an estimator $T$ of $\theta$ we measure quality by the \emph{risk}
\[
R(\theta,T)=\E_\theta\big[(T-\theta)^2\big].
\]
We would love to find the $T$ that minimises this. \textbf{But we can't}, because
$R(\theta,T)$ is a \emph{function of $\theta$}, not a number, and $\theta$ is unknown.

Concretely: whatever $T$ you propose, I beat you at the single point $\theta_0$ by proposing
the silly estimator $T'\equiv\theta_0$, which has $R(\theta_0,T')=0$. So no estimator is best
everywhere, and ``minimise the risk'' is not a well-posed problem.
\end{idea}

There are exactly three ways out of this, and the course does all three:

\begin{center}
\renewcommand{\arraystretch}{1.4}
\begin{tabular}{@{}p{3.3cm}p{5.4cm}p{6.8cm}@{}}
\toprule
\textbf{Escape route} & \textbf{What you do to $R(\theta,T)$} & \textbf{What you get} \\
\midrule
Restrict the class & Only compare estimators that are unbiased --- this rules out the silly
$T\equiv\theta_0$ & \textbf{UMVUE} (Lectures 1--5) \\
Average it & Collapse the function $R(\cdot,T)$ to a single number by averaging over $\theta$
& \textbf{Bayes} (Lectures 6--7) \\
Take the worst case & Collapse it to a single number by taking $\sup_\theta$ &
\textbf{Minimax} (Lectures 8--9) \\
\bottomrule
\end{tabular}
\end{center}

\begin{idea}{So what \emph{is} Bayes estimation, in one sentence?}
It is what you get when you turn the risk \emph{function} into a single \emph{number} by
averaging it with weights, and then minimise that number.

The weights are called the \textbf{prior}. That is the whole idea. Everything else ---
posteriors, conjugacy, Beta and Gamma distributions --- is machinery for actually doing the
computation.
\end{idea}

%%%%%%%%%%%%%%%%%%%%%%%%%%%%%%%%%%%%%%%%%%%%%%%%%%%%%%%%%
\section{The one conceptual jump: $\theta$ becomes random}

This is the step that trips everyone, so go slowly.

Everywhere in the first half of the course, $\theta$ was a \textbf{fixed unknown constant} and
$X$ was random. That is why we wrote $\E_\theta$ --- the subscript said ``the distribution
depends on this fixed number $\theta$''.

To average the risk over $\theta$, we need a weighting function on $\Theta$. Choose one and
call it $\pi(\theta)$, and require $\int_\Theta\pi(\theta)\,d\theta=1$. But a non-negative
function integrating to 1 \emph{is a probability density}. So the moment we decide to average
the risk, we have --- whether we like it or not --- treated $\theta$ as a random variable.

\begin{confusion}{``But $\theta$ isn't random, it's a fixed unknown number!''}
Correct, and the Bayesian machinery does not claim otherwise. Two readings, both legitimate:

\smallskip
\textbf{(i) Purely technical.} $\pi$ is just a set of weights you picked to turn
$R(\cdot,T)$ into a number. Nothing philosophical is being asserted. This is the reading that
makes minimax work in Lectures 8--9, where the prior is a \emph{device} for proving a theorem
about worst-case risk, and nobody believes in it.

\smallskip
\textbf{(ii) Subjective.} $\pi$ encodes what you believed about $\theta$ before seeing data,
and the posterior is your updated belief.

\smallskip
For this exam, reading (i) is enough. When a question says ``find the Bayes estimate under
prior $\pi$'', it is a computation, not a worldview.
\end{confusion}

\begin{idea}{The four objects, and how they fit together}
Once $\theta$ is random, you have a \emph{joint} distribution of the pair $(X,\theta)$, and it
can be factored two ways. Learn this table --- it is the entire subject.
\begin{center}
\renewcommand{\arraystretch}{1.5}
\begin{tabular}{@{}llp{7.6cm}@{}}
\toprule
\textbf{Name} & \textbf{Symbol} & \textbf{What it is} \\
\midrule
Prior & $\pi(\theta)$ & Distribution of $\theta$ \emph{before} seeing data. You choose it. \\
Likelihood & $f(x\mid\theta)$ & The model. Distribution of $X$ \emph{given} $\theta$ --- this
is the same $f$ you have used all semester. \\
Marginal & $m(x)=\displaystyle\int_\Theta f(x\mid\theta)\pi(\theta)\,d\theta$ &
Distribution of $X$ with $\theta$ averaged out. Usually you never compute it. \\
Posterior & $\pi(\theta\mid x)$ & Distribution of $\theta$ \emph{after} seeing $X=x$.
\textbf{This is the object you actually want.} \\
\bottomrule
\end{tabular}
\end{center}
The joint density is
\[
\underbrace{f(x\mid\theta)\,\pi(\theta)}_{\text{factor as (data given }\theta)\times\text{(prior)}}
\ =\ \text{joint}(x,\theta)\ =\
\underbrace{\pi(\theta\mid x)\,m(x)}_{\text{factor as }(\theta\text{ given data})\times\text{(marginal)}}
\]
and rearranging gives Bayes' theorem:
\[
\boxed{\ \pi(\theta\mid x)=\frac{f(x\mid\theta)\,\pi(\theta)}{m(x)}
=\frac{f(x\mid\theta)\,\pi(\theta)}{\int_\Theta f(x\mid\theta')\pi(\theta')\,d\theta'}\ }
\]
\end{idea}

\begin{idea}{The single most useful practical fact}
The denominator $m(x)$ does \emph{not} involve $\theta$. It is just whatever constant makes
$\pi(\theta\mid x)$ integrate to 1. So you almost never compute it. Instead write
\[
\pi(\theta\mid x)\ \propto\ f(x\mid\theta)\,\pi(\theta)
\qquad\text{i.e.}\qquad
\textbf{posterior}\ \propto\ \textbf{likelihood}\times\textbf{prior},
\]
throw away every factor that does not contain $\theta$, and \emph{recognise the shape} of what
is left. That is $90\%$ of the technique.
\end{idea}

%%%%%%%%%%%%%%%%%%%%%%%%%%%%%%%%%%%%%%%%%%%%%%%%%%%%%%%%%
\section{A complete worked example, slowly}

Let us do one from scratch with numbers before any theory.

\textbf{Setup.} You toss a bent coin $n=10$ times and see $x=7$ heads. So
$X\sim\mathrm{Bin}(10,\theta)$ and you want to estimate $\theta=\Prob(\text{head})$.

\textbf{Step 1: choose a prior.} Say you believe the coin is probably not wildly biased, but
you are not confident. The $\mathrm{Beta}(2,2)$ density,
$\pi(\theta)\propto\theta(1-\theta)$ on $(0,1)$, is symmetric about $\tfrac12$ and gently
humped --- a reasonable ``mildly sceptical of extremes'' choice. Its mean is $\tfrac12$.

\textbf{Step 2: write down likelihood $\times$ prior, dropping constants.}
\[
\pi(\theta\mid x=7)\ \propto\
\underbrace{\binom{10}{7}\theta^{7}(1-\theta)^{3}}_{\text{likelihood}}
\cdot
\underbrace{\frac{\theta^{2-1}(1-\theta)^{2-1}}{B(2,2)}}_{\text{prior}}
\ \propto\ \theta^{7}(1-\theta)^{3}\cdot\theta^{1}(1-\theta)^{1}
=\theta^{8}(1-\theta)^{4}.
\]
Note what got thrown away: $\binom{10}{7}$ and $1/B(2,2)$ are pure numbers, free of $\theta$.

\textbf{Step 3: recognise the shape.} A density proportional to
$\theta^{a-1}(1-\theta)^{b-1}$ is $\mathrm{Beta}(a,b)$. Here
$a-1=8$ and $b-1=4$, so
\[
\theta\mid X=7\ \sim\ \mathrm{Beta}(9,5).
\]
No integration was needed anywhere.

\textbf{Step 4: report the Bayes estimate}, which (Section 4) is the posterior mean. For
$\mathrm{Beta}(a,b)$ the mean is $a/(a+b)$:
\[
\widehat\theta_{\text{Bayes}}=\frac{9}{9+5}=\frac{9}{14}\approx0.643 .
\]

\begin{idea}{What just happened --- read this twice}
Compare three numbers:
\[
\text{prior mean}=0.500,\qquad
\widehat\theta_{\text{Bayes}}=0.643,\qquad
\text{MLE}=\frac{7}{10}=0.700 .
\]
The Bayes estimate sits \emph{between} the prior mean and the data-only estimate. That is not
a coincidence --- rearranging,
\[
\frac{a+x}{a+b+n}
=\underbrace{\frac{n}{a+b+n}}_{=10/14}\cdot\underbrace{\frac{x}{n}}_{\text{MLE}}
\;+\;\underbrace{\frac{a+b}{a+b+n}}_{=4/14}\cdot\underbrace{\frac{a}{a+b}}_{\text{prior mean}} ,
\]
and indeed $\tfrac{10}{14}(0.7)+\tfrac{4}{14}(0.5)=0.5+0.1429=0.643$. \checkmark

\smallskip
\textbf{A Bayes estimate is a weighted average of the data and the prior}, with weights
proportional to $n$ and to $a+b$. So $a+b$ behaves like a \emph{prior sample size}: a
$\mathrm{Beta}(2,2)$ prior is worth about 4 imaginary observations. As $n\to\infty$ the data
weight $\to1$ and the prior is forgotten. This picture is worth more than any formula.
\end{idea}

%%%%%%%%%%%%%%%%%%%%%%%%%%%%%%%%%%%%%%%%%%%%%%%%%%%%%%%%%
\section{Why the posterior mean? (Bayes risk)}

We have not yet justified ``take the posterior mean''. Here is where it comes from.

\begin{idea}{Definition: Bayes risk}
For a prior $\pi$ and an estimator $T$,
\[
r(\pi,T)\ :=\ \int_\Theta R(\theta,T)\,\pi(\theta)\,d\theta\ =\ \E_\pi\big[R(\theta,T)\big].
\]
\textbf{$R(\theta,T)$ is a function of $\theta$; $r(\pi,T)$ is a single number.} That is the
whole point --- numbers can be compared and minimised, functions cannot. An estimator
minimising $r(\pi,\cdot)$ is called \emph{the Bayes estimate with respect to $\pi$}, written
$T_\pi$.
\end{idea}

\begin{idea}{Theorem (the central result of Lecture 7)}
Under squared error loss, the Bayes estimate is the posterior mean:
\[
T_\pi(x)=\E[\theta\mid X=x].
\]
Moreover the minimum value achieved is
\[
r(\pi,T_\pi)=\E_m\big[\Var(\theta\mid X)\big]\qquad\text{(the average posterior variance).}
\]
\end{idea}

\textbf{Proof, with every step explained.} Start from the definition and write both
integrals out:
\[
r(\pi,T)=\int_\Theta\!\!\int_{\mathcal X}\big(T(x)-\theta\big)^2 f(x\mid\theta)\,\pi(\theta)\,dx\,d\theta .
\]
Now use the two-way factorisation of the joint density,
$f(x\mid\theta)\pi(\theta)=m(x)\,\pi(\theta\mid x)$, and swap the order of integration:
\[
r(\pi,T)=\int_{\mathcal X} m(x)\ \underbrace{\left[\int_\Theta\big(T(x)-\theta\big)^2
\pi(\theta\mid x)\,d\theta\right]}_{\text{call this }Q(x)}\,dx .
\]
This rewriting is the whole proof. Look at what it says: the Bayes risk is an average, over
$x$, of the quantity $Q(x)$, with non-negative weights $m(x)$.

So to minimise $r(\pi,T)$ it suffices to \textbf{minimise $Q(x)$ separately for each $x$}.
And for a \emph{fixed} $x$, the value $T(x)$ is just a number --- call it $a$. So we need
\[
\min_a\ \E\big[(a-\theta)^2\ \big|\ X=x\big].
\]
This is the standard fact that the mean minimises expected squared deviation:
\[
\E\big[(a-\theta)^2\mid x\big]
=\big(a-\E[\theta\mid x]\big)^2+\Var(\theta\mid x),
\]
which is minimised at $a=\E[\theta\mid x]$, with minimum value $\Var(\theta\mid x)$.
Therefore $T_\pi(x)=\E[\theta\mid X=x]$, and
$r(\pi,T_\pi)=\int m(x)\Var(\theta\mid x)dx=\E_m[\Var(\theta\mid X)]$. $\square$

\begin{idea}{Why you should care about that last formula}
$r(\pi,T_\pi)=\E_m[\Var(\theta\mid X)]$ is how \emph{every} minimax problem in this course is
computed. In the $N(\theta,1)$ problem with prior $N(0,k)$, the posterior variance is the
constant $k/(k+1)$, so the Bayes risk is $k/(k+1)$ immediately --- no integration. Memorise
it.
\end{idea}

\begin{confusion}{``Is the Bayes estimate always the posterior mean?''}
No --- \emph{only under squared error loss}. The loss function decides which feature of the
posterior you report:
\begin{center}
\begin{tabular}{@{}lll@{}}
\toprule
Loss $L(\theta,a)$ & Bayes estimate & Why \\
\midrule
$(\theta-a)^2$ (squared error) & posterior \textbf{mean} & mean minimises squared deviation \\
$|\theta-a|$ (absolute error) & posterior \textbf{median} & median minimises absolute deviation \\
$\I\{|\theta-a|>\epsilon\}$ (0--1) & posterior \textbf{mode} (MAP) & mode maximises local mass \\
\bottomrule
\end{tabular}
\end{center}
This course uses squared error throughout, so ``posterior mean'' is right here --- but say
\emph{under squared error loss} in your answer. It is a free mark.
\end{confusion}

%%%%%%%%%%%%%%%%%%%%%%%%%%%%%%%%%%%%%%%%%%%%%%%%%%%%%%%%%
\section{Conjugacy: making the computation easy}

In principle you now know everything. In practice the integral defining $m(x)$ is usually
horrible. \textbf{Conjugate priors} are the families for which it is free.

\begin{idea}{Definition}
A family $\Pi$ of priors is \textbf{conjugate} for the model $f(x\mid\theta)$ if
\[
\pi\in\Pi\quad\Longrightarrow\quad \pi(\cdot\mid x)\in\Pi\quad\text{for every }x.
\]
That is: the posterior belongs to the same family as the prior, only with updated parameters.
Updating then reduces to arithmetic on the hyperparameters.
\end{idea}

\begin{recipe}{How to find or verify a conjugate family --- the kernel-matching drill}
\begin{enumerate}
\item Write the likelihood and \textbf{delete every factor free of $\theta$}. What survives is
the \emph{kernel}.
\item Ask: what family of $\theta$-densities has this same functional shape? That family is
the conjugate one.
\item Multiply prior kernel $\times$ likelihood kernel and collect exponents.
\item Read off the updated parameters. Never integrate.
\end{enumerate}
\end{recipe}

\subsection{Beta--Binomial}
$X\sim\mathrm{Bin}(n,\theta)$. Likelihood kernel: $\theta^{x}(1-\theta)^{n-x}$ --- a power of
$\theta$ times a power of $(1-\theta)$. Which densities on $(0,1)$ look like that?
$\mathrm{Beta}(a,b)\propto\theta^{a-1}(1-\theta)^{b-1}$. Multiply:
\[
\theta^{x}(1-\theta)^{n-x}\cdot\theta^{a-1}(1-\theta)^{b-1}=\theta^{(a+x)-1}(1-\theta)^{(b+n-x)-1}
\]
\[
\boxed{\ \theta\mid x\sim\mathrm{Beta}(a+x,\ b+n-x),\qquad
T_\pi=\frac{a+X}{a+b+n}\ }
\]
\emph{Interpretation:} add the successes to $a$, the failures to $b$.

\subsection{Gamma--Poisson}
$X_1,\dots,X_n\iid\mathrm{Poi}(\lambda)$, $S=\sum X_i$. Likelihood kernel:
$e^{-n\lambda}\lambda^{S}$ --- a power of $\lambda$ times $e^{-\text{const}\cdot\lambda}$.
Which densities on $(0,\infty)$ look like that? $\mathrm{Gamma}(\alpha,\text{scale }\beta)
\propto\lambda^{\alpha-1}e^{-\lambda/\beta}$. Multiply:
\[
e^{-n\lambda}\lambda^{S}\cdot\lambda^{\alpha-1}e^{-\lambda/\beta}
=\lambda^{(\alpha+S)-1}\exp\Big\{-\lambda\Big(n+\frac1\beta\Big)\Big\}
\]
\[
\boxed{\ \lambda\mid\mathbf x\sim\mathrm{Gamma}\Big(\alpha+S,\ \frac{\beta}{1+n\beta}\Big),
\qquad T_\pi=\frac{\beta(\alpha+S)}{1+n\beta}\ }
\]
\emph{Numerical check.} $n=5$, $S=12$, prior $\mathrm{Gamma}(2,1)$ (mean 2). Posterior is
$\mathrm{Gamma}(14,\tfrac16)$ with mean $14/6\approx2.333$. Compare MLE $=12/5=2.4$ and prior
mean $=2$: again strictly between. \checkmark

\subsection{Normal--Normal}
$X_1,\dots,X_n\iid N(\theta,\sigma^2)$ with $\sigma^2$ \emph{known}; prior $\theta\sim N(\mu,\tau^2)$.
The exponent is a sum of two quadratics in $\theta$, so completing the square gives a normal
again. The cleanest way to state the answer uses \textbf{precision} $=1/\text{variance}$:
\[
\boxed{\
\text{posterior precision}=\frac{n}{\sigma^{2}}+\frac{1}{\tau^{2}},\qquad
\text{posterior mean}=
\frac{\dfrac{n}{\sigma^{2}}\Xbar+\dfrac{1}{\tau^{2}}\mu}
     {\dfrac{n}{\sigma^{2}}+\dfrac{1}{\tau^{2}}}\ }
\]
In words: \emph{precisions add, and the posterior mean is the precision-weighted average of
the data mean and the prior mean.} If you remember only one thing about Normal--Normal,
remember that sentence --- you can reconstruct the algebra from it.

\emph{Numerical check.} $n=4$, $\Xbar=10$, $\sigma^2=4$ (so $\sigma^2/n=1$, data precision 1);
prior $N(8,4)$ (precision $\tfrac14$). Posterior precision $=1+\tfrac14=1.25$, so posterior
variance $=0.8$; posterior mean $=\dfrac{1\cdot10+\tfrac14\cdot8}{1.25}=\dfrac{12}{1.25}=9.6$.
Between the prior mean 8 and the data mean 10, closer to the data because the data are four
times as precise. \checkmark

\subsection{Summary table}
\begin{center}
\renewcommand{\arraystretch}{1.5}
\begin{tabular}{@{}llll@{}}
\toprule
Model & Conjugate prior & Posterior & Bayes estimate \\
\midrule
$\mathrm{Bin}(n,\theta)$ & $\mathrm{Beta}(a,b)$ & $\mathrm{Beta}(a+x,\,b+n-x)$ & $\dfrac{a+X}{a+b+n}$ \\
$\mathrm{Poi}(\lambda)$, $n$ obs & $\mathrm{Gamma}(\alpha,\beta)$ & $\mathrm{Gamma}\big(\alpha+S,\tfrac{\beta}{1+n\beta}\big)$ & $\dfrac{\beta(\alpha+S)}{1+n\beta}$ \\
$N(\theta,\sigma^2)$, $n$ obs & $N(\mu,\tau^2)$ & $N\big(\text{prec.\ wtd.\ mean},\ (\tfrac{n}{\sigma^2}+\tfrac1{\tau^2})^{-1}\big)$ & precision-weighted mean \\
$\mathrm{Exp}(\theta)$, $n$ obs & Inverse-Gamma$(\alpha,\beta)$ & Inv-Gamma$(\alpha+n,\ \beta+S)$ & $\dfrac{\beta+S}{\alpha+n-1}$ \\
$U[a,\theta]$, $n$ obs & $\mathrm{Pareto}(\alpha,a)$ & $\mathrm{Pareto}\big(\alpha+n,\max(a,X_{(n)})\big)$ & $\dfrac{(\alpha+n)M}{\alpha+n-1}$ \\
\bottomrule
\end{tabular}
\end{center}

\begin{confusion}{``Do I have to memorise this table?''}
No, and you should not try. Every row is 30 seconds of kernel matching. What you \emph{must}
know cold is the \emph{drill}: strip constants, match the shape, add the exponents. In the
exam, derive the row you need --- it is more reliable than a half-remembered formula, and
derivations earn marks that quoted answers do not.
\end{confusion}

\begin{confusion}{``Why can't the Cauchy have a nice conjugate family?''}
Because conjugacy is really a property of \emph{exponential families}. In an exponential
family, $\prod_i f(x_i\mid\theta)$ depends on the data only through $\sum_iT(x_i)$ and $n$,
so a two-parameter prior family suffices --- and the update is just ``add $\sum T(x_i)$ to one
hyperparameter and $n$ to the other''. Look back at the table: every update has exactly that
form.

The Cauchy is not an exponential family, so its likelihood cannot be compressed into a fixed
number of summaries. A conjugate family still exists (products of Cauchy kernels), but it has
to grow by one hyperparameter per observation --- useless in practice. This is the point of
the ``find a conjugate family for any density'' exercise in Lecture 7: the class of \emph{all}
densities is trivially conjugate, and equally useless. \textbf{Conjugacy is only valuable
when the family is small.}
\end{confusion}

%%%%%%%%%%%%%%%%%%%%%%%%%%%%%%%%%%%%%%%%%%%%%%%%%%%%%%%%%
\section{Three properties you will be asked about}

\subsection{A Bayes estimate is essentially never unbiased}

You saw it numerically in Section 3: the estimate is pulled toward the prior mean. That pull
is \emph{shrinkage}, and shrinkage is exactly bias. Concretely, for Beta--Binomial,
\[
\E_\theta\Big[\frac{a+X}{a+b+n}\Big]=\frac{a+n\theta}{a+b+n},
\]
which equals $\theta$ for all $\theta$ only if $a=0$ and $a+b=0$ --- not a proper prior.

The general argument is four lines and is a standard exam question. Suppose $T$ is both
Bayes ($T=\E[\theta\mid X]$) and unbiased ($\E[T\mid\theta]=\theta$). Write $\E_m$ for
expectation under the joint law and condition two different ways:
\[
\E_m[T\theta]=\E\big[\theta\,\E(T\mid\theta)\big]=\E[\theta^2],
\qquad
\E_m[T\theta]=\E\big[T\,\E(\theta\mid X)\big]=\E[T^2].
\]
Then
\[
\E_m\big[(T-\theta)^2\big]=\E[T^2]+\E[\theta^2]-2\E_m[T\theta]=0,
\]
which forces $T=\theta$ almost surely --- impossible unless the problem is degenerate.

\begin{confusion}{``So Bayes estimates are worse than UMVUEs?''}
No. Bias is not a defect, it is a \emph{trade}. Shrinking toward the prior costs you a little
bias and buys you a lot of variance reduction, and MSE is the sum of the two. You already met
this in Lecture 1: $\frac{1}{n+1}\sum(X_i-\Xbar)^2$ is biased and beats the unbiased $S^2$ at
every $\sigma^2$. And in Lecture 8 the \emph{minimax} estimator of a binomial proportion is a
biased Bayes estimate, not the UMVUE. Unbiasedness is a constraint we impose for
tractability, not a virtue.
\end{confusion}

\subsection{The posterior depends on the data only through a sufficient statistic}

If $f(\mathbf x\mid\theta)=g(\theta,T(\mathbf x))\,h(\mathbf x)$ by the factorisation theorem,
then
\[
\pi(\theta\mid\mathbf x)=\frac{g(\theta,T(\mathbf x))\,h(\mathbf x)\,\pi(\theta)}
{\int g(\theta',T(\mathbf x))\,h(\mathbf x)\,\pi(\theta')\,d\theta'}
=\frac{g(\theta,T(\mathbf x))\,\pi(\theta)}{\int g(\theta',T(\mathbf x))\,\pi(\theta')\,d\theta'} ,
\]
since $h(\mathbf x)$ cancels between numerator and denominator. So the posterior --- and hence
every Bayes estimate --- is a function of $T$ alone. Check the table: every entry depends on
the data only through $\sum X_i$, $\Xbar$, or $X_{(n)}$. \checkmark

\subsection{Improper priors and generalized Bayes}

Sometimes you want to express ``no prior information''. The natural choice on $\theta\in\R$ is
$\pi(\theta)\equiv1$ --- but $\int_\R 1\,d\theta=\infty$, so this is \textbf{not} a probability
density. Such a $\pi$ is called \emph{improper}.

Remarkably, the posterior can still be a perfectly good density. If it is, its mean is called
the \textbf{generalized Bayes estimate}.

\begin{idea}{The three examples from the course}
\begin{enumerate}
\item $X_1,\dots,X_n\iid N(\theta,1)$, $\pi\equiv1$. Then
$\pi(\theta\mid\mathbf x)\propto\exp\{-\tfrac n2(\theta-\bar x)^2\}$, i.e.\
$N(\bar x,1/n)$ --- proper. Generalized Bayes estimate: $\boxed{\Xbar}$. Note this is also the
MLE and the UMVUE. (Sanity check on Normal--Normal: let $\tau^2\to\infty$, so prior precision
$\to0$ and the posterior mean $\to\Xbar$. \checkmark)
\item $X_1,\dots,X_n\iid U[\theta-\tfrac12,\theta+\tfrac12]$, $\pi\equiv1$. The likelihood is
the indicator $\I[x_{(n)}-\tfrac12\le\theta\le x_{(1)}+\tfrac12]$, so the posterior is
\emph{uniform} on that interval and the estimate is its midpoint:
$\boxed{\tfrac12(X_{(1)}+X_{(n)})}$, the midrange. (This is Homework 2 Q3.)
\item $X_1,\dots,X_n\iid\mathrm{Exp}(\theta)$, $\pi(\theta)=1/\theta$. Posterior is
Inverse-Gamma$(n,S)$ with mean $\boxed{S/(n-1)}$ for $n\ge2$ --- note this is \emph{not}
$\Xbar$.
\end{enumerate}
\end{idea}

\begin{confusion}{``If the prior is improper, is the estimate still Bayes?''}
Not in the technical sense --- and this matters in Lecture 8. The theorem ``Bayes with
constant risk $\Rightarrow$ minimax'' needs $\pi$ to be a \emph{probability} measure, because
its proof uses ``$\sup\ \ge$ average''. Under an improper prior the Bayes risk is infinite and
that step is meaningless.

That is exactly why $X$ being minimax for $N(\theta,1)$ needs the \emph{sequence} argument
with proper priors $\pi_k=N(0,k)$, rather than the one-line constant-risk argument. If you
are asked ``can a generalized Bayes equalizer estimate be minimax?'', the answer is: yes it
can, but you must prove it via limits of proper Bayes rules, not via the constant-risk
theorem.
\end{confusion}

%%%%%%%%%%%%%%%%%%%%%%%%%%%%%%%%%%%%%%%%%%%%%%%%%%%%%%%%%
\section{How Bayes connects to the rest of the course}

\begin{idea}{The map}
\[
\textbf{Bayes}\ \xrightarrow[\text{take }\sup_\theta]{\text{constant risk}}\ \textbf{Minimax}
\]
The reason Bayes appears immediately before minimax is that \emph{Bayes estimates are the tool
for proving minimaxity}. Two theorems:
\begin{enumerate}
\item \textbf{Constant-risk criterion.} If $T_\pi$ is Bayes for a proper $\pi$ and
$R(\theta,T_\pi)=c$ for all $\theta$, then $T_\pi$ is minimax. Proof:
$\sup_\theta R(\theta,T)\ge r(\pi,T)\ge r(\pi,T_\pi)=c=\sup_\theta R(\theta,T_\pi)$.
\item \textbf{Limit of Bayes rules.} If $r(\pi_k,T_{\pi_k})\to c$ and
$\sup_\theta R(\theta,T^*)=c$, then $T^*$ is minimax.
\end{enumerate}
In both, the middle inequality ``$\sup\ \ge$ average'' is what does the work --- and averaging
is precisely what the prior was introduced for in Section 1. The circle closes.
\end{idea}

%%%%%%%%%%%%%%%%%%%%%%%%%%%%%%%%%%%%%%%%%%%%%%%%%%%%%%%%%
\section{Answering an exam question}

\begin{recipe}{The standard six lines}
\begin{enumerate}
\item Write \textbf{posterior $\propto$ likelihood $\times$ prior}.
\item Delete every factor free of $\theta$.
\item Collect exponents; \textbf{name} the resulting standard density.
\item Quote its mean. State ``\emph{the Bayes estimate under squared error loss is the
posterior mean}''.
\item If asked about bias, compute $\E_\theta T$ directly and observe it is not $\theta$.
\item If asked for the Bayes \emph{risk}, use $r(\pi,T_\pi)=\E_m[\Var(\theta\mid X)]$ --- take
the posterior variance and average it over the marginal.
\end{enumerate}
\end{recipe}

%%%%%%%%%%%%%%%%%%%%%%%%%%%%%%%%%%%%%%%%%%%%%%%%%%%%%%%%%
\section{Practice, with answers}

Do these with the answers covered. They take about 25 minutes together.

\begin{enumerate}[label=\textbf{P\arabic*.}]
\item $X_1,\dots,X_n\iid\mathrm{Bern}(\theta)$ with prior $\mathrm{Beta}(a,b)$. Find the Bayes
estimate, express it as a weighted average of the MLE and the prior mean, and say what happens
as $n\to\infty$.

\item $X\sim\mathrm{Geometric}$, $f(x\mid\theta)=\theta(1-\theta)^{x}$, $x=0,1,2,\dots$
Find a conjugate family and the Bayes estimate of $\theta$.

\item $X_1,\dots,X_n\iid N(\theta,\sigma^2)$, $\sigma^2$ known, prior $N(\mu,\tau^2)$. Show
the posterior variance is strictly smaller than both $\sigma^2/n$ and $\tau^2$. Interpret.

\item $X_1,\dots,X_n\iid U[0,\theta]$ with prior $\mathrm{Pareto}(\alpha,m)$, i.e.\
$\pi(\theta)=\alpha m^{\alpha}/\theta^{\alpha+1}$ for $\theta>m$. Find the posterior and the
Bayes estimate.

\item $X\sim N(\theta,1)$ with prior $N(0,k)$. Find the Bayes estimate and show the Bayes
risk equals $k/(k+1)$. What happens as $k\to\infty$, and why does it matter?
\end{enumerate}

\subsection*{Answers}

\noindent\textbf{P1.} Posterior $\propto\theta^{S}(1-\theta)^{n-S}\theta^{a-1}(1-\theta)^{b-1}$
with $S=\sum X_i$, so $\mathrm{Beta}(a+S,\ b+n-S)$ and
\[
T_\pi=\frac{a+S}{a+b+n}
=\frac{n}{a+b+n}\cdot\Xbar+\frac{a+b}{a+b+n}\cdot\frac{a}{a+b}.
\]
As $n\to\infty$ the data weight $\to1$, so $T_\pi-\Xbar\to0$: \textbf{the prior washes out}.
This is why the choice of prior matters little for large samples and a lot for small ones.

\smallskip
\noindent\textbf{P2.} Likelihood kernel $\theta(1-\theta)^{x}$ --- a power of $\theta$ times a
power of $(1-\theta)$, so try $\mathrm{Beta}(a,b)$ again:
\[
\theta(1-\theta)^{x}\cdot\theta^{a-1}(1-\theta)^{b-1}=\theta^{(a+1)-1}(1-\theta)^{(b+x)-1},
\]
so $\theta\mid x\sim\mathrm{Beta}(a+1,\ b+x)$ and
$T_\pi=\dfrac{a+1}{a+b+x+1}$. The Beta family is conjugate for the geometric too --- as it
must be, since both are exponential families in $\theta$ with the same $(1-\theta)$ structure.

\smallskip
\noindent\textbf{P3.} Posterior variance
$=\Big(\dfrac{n}{\sigma^2}+\dfrac{1}{\tau^2}\Big)^{-1}=\dfrac{\sigma^2\tau^2}{n\tau^2+\sigma^2}$.
Then
\[
\frac{\sigma^2\tau^2}{n\tau^2+\sigma^2}<\frac{\sigma^2}{n}
\iff n\tau^2<n\tau^2+\sigma^2\ \checkmark,
\qquad
\frac{\sigma^2\tau^2}{n\tau^2+\sigma^2}<\tau^2
\iff \sigma^2<n\tau^2+\sigma^2\ \checkmark .
\]
\emph{Interpretation:} precisions add, so combining two sources of information is always at
least as good as either alone. The posterior is sharper than the data alone \emph{and} sharper
than the prior alone.

\smallskip
\noindent\textbf{P4.} Likelihood $\theta^{-n}\I[\theta\ge x_{(n)}]$, prior
$\propto\theta^{-(\alpha+1)}\I[\theta>m]$. Multiply:
\[
\pi(\theta\mid\mathbf x)\propto\theta^{-(\alpha+n)-1}\,\I\big[\theta>\max(m,x_{(n)})\big]
\]
which is $\mathrm{Pareto}\big(\alpha+n,\ M\big)$ with $M=\max(m,X_{(n)})$. A
$\mathrm{Pareto}(\beta,M)$ has mean $\beta M/(\beta-1)$ for $\beta>1$, so
\[
T_\pi=\frac{(\alpha+n)\,M}{\alpha+n-1}.
\]
Note the shape: like the UMVUE $\tfrac{n+1}{n}X_{(n)}$, it inflates the largest observation ---
sensible, since $X_{(n)}$ always \emph{under}estimates $\theta$.

\smallskip
\noindent\textbf{P5.} Precision form: data precision 1, prior precision $1/k$, so posterior
precision $=1+\tfrac1k=\tfrac{k+1}{k}$ and
\[
\theta\mid X=x\ \sim\ N\Big(\frac{k}{k+1}x,\ \frac{k}{k+1}\Big),
\qquad T_\pi=\frac{k}{k+1}X .
\]
The posterior variance is the \emph{constant} $k/(k+1)$ (it does not involve $x$), so
\[
r(\pi_k,T_{\pi_k})=\E_m\big[\Var(\theta\mid X)\big]=\frac{k}{k+1}.
\]
As $k\to\infty$ this $\to1$, which is exactly $\sup_\theta R(\theta,X)=1$. By the
limit-of-Bayes-rules theorem, $X$ is \textbf{minimax}. This is Homework 2 Q1, and it is the
single most important place Bayes machinery is used in this course --- note that $X$ itself is
not Bayes for any proper prior; the sequence is what does the work.

\vfill
\begin{center}
\small\itshape
Companions: \texttt{PI\_Exam\_Revision.pdf} \S8--\S9,
\texttt{PI\_Class\_Exercises\_Solved.pdf} \S6--\S7,
\texttt{PI\_Mock\_Exam.pdf} Q5 and Q6.
\end{center}

\end{document}
